\section{A tétel lényege és vizsgaszintű vázlata}

\begin{summarybox}
A 25. tétel a szoftvertechnológia és a szoftverfolyamat alapfogalmait foglalja össze. A szoftvertechnológia lényege, hogy a szoftvereket ne alkalmi, egyéni programozási munkaként, hanem mérnöki szemlélettel, tervezetten, dokumentáltan, költség- és határidőkorlátokat figyelembe véve hozzuk létre és tartsuk karban. A szoftverfolyamat azoknak a tevékenységeknek és eredményeknek a rendszere, amelyek egy szoftvertermék előállításához vezetnek. A közös folyamatfázisok: specifikáció, tervezés, implementáció, validáció és evolúció. A tételben különösen fontos az evolúciós modell, amely működő verziók és felhasználói visszajelzések segítségével fokozatosan alakítja ki a végleges rendszert.
\end{summarybox}

Vizsgán célszerű a következő sorrendet követni:

\begin{enumerate}
  \item szoftvertechnológia fogalma és célja;
  \item szoftvertermékek rövid osztályozása: általános termék és egyedi igényre szabott termék;
  \item szoftverfolyamat fogalma, emberi és szervezeti jellege;
  \item kiegészítő munkafolyamatok: projektmenedzsment, verziókezelés, erőforrás-kezelés, minőségbiztosítás, terméktámogatás;
  \item a folyamat közös fázisai: szoftverspecifikáció, tervezés, implementáció, validáció, evolúció;
  \item szoftverfolyamat-modellek fogalma és fontosabb példák;
  \item vízesésmodell rövid szerepe, előnye és korlátja;
  \item evolúciós fejlesztés alapgondolata;
  \item feltáró fejlesztés és eldobható prototípuskészítés összehasonlítása;
  \item evolúciós modell előnyei, hátrányai és alkalmazhatósága;
  \item rövid összevetés inkrementális, komponensalapú és spirális szemlélettel;
  \item gyakori vizsgahibák és rövid összefoglaló válasz.
\end{enumerate}

A legfontosabb vizsgamondat: a szoftverfolyamat nem pusztán programozás, hanem a követelmények megértésétől a tervezésen, implementáción és validáción át a szoftver teljes életciklusa alatt végzett evolúcióig tartó, szervezett tevékenységrendszer; az evolúciós modell ezt nem merev lineáris sorrendként, hanem működő verziók és visszacsatolások sorozataként kezeli.

\section{A szoftvertechnológia fogalma}

\subsection{Miért van szükség szoftvertechnológiára?}

A kis, egy ember által írt programoknál gyakran elegendő, ha a fejlesztő fejben tartja a feladatot, gyorsan leírja a megoldást, lefuttatja, majd javítja a hibákat. Nagyobb szoftverrendszereknél ez már nem működik megbízhatóan. A rendszer több ember munkája, hosszabb ideig készül, sok követelménynek kell megfelelnie, dokumentációt igényel, később üzemeltetni és módosítani kell.

A \textbf{szoftvertechnológia} olyan mérnöki, technológiai és vezetési alapelvek összessége, amelyek a szoftverek termékszerű előállítását, használatát és karbantartását támogatják. Nem csak programozási technikákat jelent, hanem módszereket a követelmények feltárására, a rendszer megtervezésére, a minőség ellenőrzésére, a dokumentáció elkészítésére, a projekt vezetésére és a változások kezelésére.

Másképpen megfogalmazva: a szoftvertechnológia célja, hogy a szoftverfejlesztés ne véletlenszerű próbálkozások sorozata legyen, hanem tervezhető, ellenőrizhető, ismételhető és javítható folyamat. A cél nem az, hogy minden projekt ugyanúgy készüljön, hanem az, hogy az adott szervezet, csapat, rendszer és megrendelő sajátosságaihoz illő folyamatot alakítsunk ki.

\subsection{Szoftvertermékek csoportjai}

A szoftvertermékeket gyakran két nagy csoportba soroljuk.

\textbf{Általános termékek} esetén a szoftvert egy fejlesztő szervezet készíti, majd több ügyfélnek, felhasználónak értékesíti. Ilyenek például a szövegszerkesztők, adatbázis-kezelők, fejlesztőeszközök vagy általános irodai alkalmazások. Itt a fejlesztő cég dönt a funkciókról, a piac igényeit és a versenytársakat figyelembe véve.

\textbf{Egyedi igényekre szabott termékek} esetén a szoftver konkrét megrendelői igény alapján készül. Ilyenek lehetnek ipari vezérlőrendszerek, forgalomirányító rendszerek, vállalati belső rendszerek vagy egy szervezet saját munkafolyamatait támogató alkalmazások. Ilyenkor a megrendelő követelményei, szabályai és környezete közvetlenül meghatározzák a fejlesztést.

A határ nem mindig éles. Egy általános termék is testreszabható, és egy egyedi rendszer is épülhet kész komponensekre. A szoftvertechnológiai szemlélet mindkét esetben fontos, mert mindkét esetben kezelni kell a követelményeket, a minőséget, a változásokat és a karbantartást.

\section{A szoftverfolyamat fogalma}

\subsection{A szoftverfolyamat jelentése}

A \textbf{szoftverfolyamat} tevékenységek és eredmények sora, amelyek egy szoftvertermék előállításához vezetnek. Tevékenység például a követelmények feltárása, a tervezés, a programozás, a tesztelés, a dokumentálás vagy a verzió kiadása. Eredmény lehet követelménydokumentum, tervdokumentum, forráskód, tesztjelentés, futtatható rendszer vagy új szoftververzió.

A szoftverfolyamat összetett és erősen függ az emberi tevékenységektől. Emiatt nem lehet teljesen automatizálni. Eszközökkel lehet támogatni, például verziókezelővel, tesztkeretrendszerrel, modellezőeszközzel vagy projektmenedzsment-rendszerrel, de a lényegi döntések jelentős része továbbra is emberi: mit kér a megrendelő, mi a jó architektúra, mennyi kockázat vállalható, mikor tekinthető elfogadhatónak a rendszer.

Nincs egyetlen ideális, minden projektre megfelelő folyamat. Más folyamat illik egy kis webalkalmazáshoz, mint egy biztonságkritikus repülésirányító rendszerhez. A jó szoftverfolyamat figyelembe veszi a szervezet képességeit, a fejlesztők tapasztalatát, az eszközöket, az ügyfél bevonhatóságát, a rendszer méretét és a minőségi követelményeket.

\subsection{Kiegészítő munkafolyamatok}

A fő fejlesztési tevékenységek mellett több olyan kiegészítő munkafolyamat van, amely nélkül a nagyobb projektek nehezen lennének kezelhetők.

\begin{itemize}
  \item \textbf{Projektmenedzsment}: feladatok, határidők, költségek, felelősségek és erőforrások tervezése, követése.
  \item \textbf{Verziókezelés és konfigurációkezelés}: forráskódok, dokumentumok, buildelt változatok és kiadások követése.
  \item \textbf{Erőforrás-menedzsment}: emberek, eszközök, infrastruktúra és időkeretek beosztása.
  \item \textbf{Minőségbiztosítás}: módszerek és ellenőrzések annak biztosítására, hogy a termék és a folyamat megfelelő minőségű legyen.
  \item \textbf{Terméktámogatás}: a már átadott rendszer használatának támogatása, hibák fogadása, felhasználói problémák kezelése.
  \item \textbf{Folyamatértékelés és továbbfejlesztés}: annak vizsgálata, hogy a használt fejlesztési módszer mennyire működött jól, és mit kell rajta javítani.
\end{itemize}

Ezek a tevékenységek nem mellékesek. Egy projekt technikailag jó kódot is előállíthat, de sikertelen lehet, ha közben rosszul kezelik a követelményeket, nincs verziókövetés, nincs tesztelési stratégia, vagy a módosítások követhetetlenné válnak.

\section{A szoftverfolyamat közös fázisai}

\subsection{Áttekintés}

A legtöbb szoftverfolyamatban felismerhető néhány közös tevékenység. A különböző modellek abban térnek el, hogy ezeket a tevékenységeket milyen sorrendben, mennyire szigorúan, mennyi visszacsatolással és milyen dokumentáltsággal végzik.

\begin{center}
\begin{tabularx}{\textwidth}{L{31mm}X X}
\toprule
\textbf{Fázis} & \textbf{Lényege} & \textbf{Jellemző eredménye} \\
\midrule
Szoftverspecifikáció & A rendszer szolgáltatásainak, korlátainak, céljainak és minőségi elvárásainak megértése és rögzítése. & Követelménydokumentum, felhasználói és rendszerkövetelmények, megvalósíthatósági jelentés. \\
Tervezés & A specifikáció alapján kialakul a rendszer szerkezete: architektúra, alrendszerek, interfészek, adatszerkezetek és algoritmusok. & Szoftverterv, architekturális modell, interfészleírások, komponens- és adatmodell. \\
Implementáció & A terv futtatható programmá alakítása; a komponensek elkészítése és fejlesztői szintű ellenőrzése. & Forráskód, programegységek, buildelhető rendszer, egységtesztek. \\
Validáció & Annak ellenőrzése, hogy a rendszer megfelel-e a specifikációnak és a megrendelő valódi igényeinek. & Teszttervek, tesztjegyzőkönyvek, hibajegyek, elfogadási eredmény. \\
Evolúció & A már átadott vagy használatban lévő rendszer módosítása hibák, új követelmények és környezeti változások miatt. & Javítások, új verziók, karbantartási és továbbfejlesztési döntések. \\
\bottomrule
\end{tabularx}

\end{center}

\begin{figure}[htbp]
\centering
\resizebox{0.96\textwidth}{!}{%
\begin{tikzpicture}[font=\scriptsize, node distance=7mm]
  \tikzset{phase/.style={draw, rounded corners, align=center, minimum width=29mm, minimum height=11mm},
           note/.style={draw, rounded corners, align=center, fill=black!5, minimum width=36mm, minimum height=8mm}}

  \node[phase, fill=warnbg] (spec) {Szoftver-\\specifikáció};
  \node[phase, fill=lightaccent, right=of spec] (design) {Tervezés};
  \node[phase, fill=lightaccent, right=of design] (impl) {Implementáció};
  \node[phase, fill=warnbg, right=of impl] (val) {Validáció};
  \node[phase, fill=black!5, right=of val] (evo) {Evolúció};

  \draw[arrow] (spec) -- (design);
  \draw[arrow] (design) -- (impl);
  \draw[arrow] (impl) -- (val);
  \draw[arrow] (val) -- (evo);

  \draw[arrow, dashed] ([yshift=-2mm]val.south) .. controls +(0,-13mm) and +(0,-13mm) .. ([yshift=-2mm]spec.south);
  \draw[arrow, dashed] ([yshift=2mm]evo.north) .. controls +(0,15mm) and +(0,15mm) .. node[above, align=center]{új verziók és változtatások} ([yshift=2mm]design.north);

  \node[note, below=13mm of impl] {A gyakorlatban a fázisok nem mindig tisztán egymás után következnek; a modellek azt adják meg, hogyan szervezzük őket.};
\end{tikzpicture}%
}
\caption{A szoftverfolyamat közös fázisai és a visszacsatolások szerepe}
\label{fig:zv25_szoftverfolyamat_fazisok}
\end{figure}


\subsection{Szoftverspecifikáció}

A \textbf{szoftverspecifikáció}, más néven követelménytervezés az a folyamat, amelyben megértjük és definiáljuk, hogy a rendszernek milyen szolgáltatásokat kell biztosítania, milyen megszorítások között kell működnie, és milyen minőségi elvárásoknak kell megfelelnie.

Ez kritikus szakasz, mert a specifikációban vétett hibák később sokkal drágábban javíthatók. Ha rosszul értelmezzük a felhasználó igényét, akkor a tervezés és az implementáció akár hibátlanul is végrehajtható, mégsem azt a rendszert kapjuk, amire szükség lenne.

A szoftverspecifikáció gyakori résztevékenységei:

\begin{enumerate}
  \item \textbf{Megvalósíthatósági tanulmány}: eldönti, hogy a kívánt rendszer technikailag és gazdaságilag kivitelezhető-e.
  \item \textbf{Követelmények feltárása és elemzése}: meglévő rendszerek vizsgálata, interjúk, megbeszélések, megfigyelések, prototípusok és modellek segítségével összegyűjti az igényeket.
  \item \textbf{Követelményspecifikáció}: az összegyűjtött információkat egységes dokumentummá alakítja. Megkülönböztethetünk felhasználói követelményeket és részletes rendszerkövetelményeket.
  \item \textbf{Követelményvalidáció}: ellenőrzi, hogy a követelmények valószerűek, konzisztensek, teljesek és érthetők-e.
\end{enumerate}

Fontos, hogy ezek a lépések nem mindig szigorúan egymás után következnek. Gyakori, hogy egy feltárt követelmény újabb elemzést igényel, vagy validáció közben derül ki, hogy egy korábbi követelmény hiányos vagy ellentmondásos.

\subsection{Szoftvertervezés}

A \textbf{szoftvertervezés} a specifikáció futtatható rendszer felé történő közelítése. A tervezésben azt döntjük el, hogyan épüljön fel a rendszer, milyen fő részekből álljon, ezek hogyan kommunikáljanak, milyen adatokat tároljanak, és milyen algoritmusokat használjanak.

A tervezés tipikus lépései:

\begin{itemize}
  \item \textbf{architekturális tervezés}: alrendszerek és kapcsolataik azonosítása;
  \item \textbf{absztrakt specifikáció}: az alrendszerek szolgáltatásainak és megszorításainak megadása;
  \item \textbf{interfésztervezés}: az alrendszerek és komponensek közötti kapcsolódási pontok meghatározása;
  \item \textbf{komponenstervezés}: szolgáltatások elhelyezése komponensekben;
  \item \textbf{adatszerkezet-tervezés}: az adatok belső reprezentációjának és tárolásának megtervezése;
  \item \textbf{algoritmustervezés}: a szolgáltatásokhoz szükséges algoritmusok meghatározása.
\end{itemize}

A tervezés iteratív tevékenység. A részletek kidolgozása közben gyakran derülnek ki korábbi félreértések vagy hiányosságok. Ilyenkor vissza kell csatolni a specifikációhoz, módosítani kell a tervet, vagy új döntést kell hozni az architektúráról.

\subsection{Implementáció}

Az \textbf{implementáció} során a tervből forráskód, majd futtatható rendszer lesz. Ide tartozik a komponensek programozása, a buildelés, a fejlesztői tesztelés, a hibajavítás és a kód integrálása.

A programozás erősen személyekhez és technológiákhoz kötött tevékenység, ezért nincs minden helyzetre érvényes általános szabály. A gyakorlatban gyakran azokkal a komponensekkel kezdik a fejlesztést, amelyeket a csapat a legjobban ért, vagy amelyek magas kockázatúak, és ezért korai bizonyítást igényelnek. Máskor a legfontosabb üzleti funkciók kapnak elsőbbséget.

Az implementáció nem választható el teljesen a teszteléstől. A fejlesztők már programozás közben is ellenőrzik a komponenseket, futtatják az egységteszteket, kipróbálják az interfészeket, és hibák esetén javítanak. Nagyobb rendszerben a komponensek külön-külön működhetnek, de integráció után mégis hibák jelenhetnek meg az együttműködésükben.

\subsection{Validáció, verifikáció és tesztelés}

A \textbf{validáció} célja annak ellenőrzése, hogy a megfelelő terméket készítettük-e el, vagyis a rendszer megfelel-e a felhasználó és a megrendelő valódi igényeinek. A \textbf{verifikáció} azt vizsgálja, hogy a terméket jól készítettük-e el, vagyis megfelel-e a specifikációnak.

Rövid vizsgamondatként megjegyezhető:

\begin{itemize}
  \item \textbf{verifikáció}: jól készítjük-e a terméket;
  \item \textbf{validáció}: a megfelelő terméket készítjük-e.
\end{itemize}

A validáció és verifikáció a fejlesztés minden lépésében jelen lehet. Statikus technika például a követelménydokumentum, tervdokumentum vagy forráskód átvizsgálása. Dinamikus technika a tesztelés, ahol a programot konkrét tesztadatokkal futtatjuk.

A tesztelés gyakori szintjei:

\begin{enumerate}
  \item \textbf{komponens- vagy egységtesztelés}: egyedi komponensek ellenőrzése;
  \item \textbf{rendszertesztelés}: az alrendszerek és komponensek együttműködésének, interfészeinek és rendszerszintű tulajdonságainak vizsgálata;
  \item \textbf{átvételi tesztelés}: a megrendelő vagy végfelhasználó adataival és szempontjai szerint ellenőrzi, hogy a rendszer elfogadható-e.
\end{enumerate}

Ha a validáció hibát tár fel, akkor nemcsak a kódot kell javítani. Előfordulhat, hogy a specifikáció, a terv, a tesztek vagy akár az egész fejlesztési stratégia módosításra szorul.

\subsection{Szoftverevolúció}

A \textbf{szoftverevolúció} a szoftver változtatásának folyamata a rendszer életciklusa során. Régebbi szemléletben élesebben elválasztották a fejlesztést és a karbantartást: előbb elkészült a rendszer, majd átadás után következett a karbantartás. A modern szemléletben ez a határ sokkal kevésbé éles, mert kevés rendszer tekinthető teljesen újnak, és a legtöbb szoftver folyamatosan változik.

A szoftverevolúció oka lehet:

\begin{itemize}
  \item később kiderülő hibák javítása;
  \item új felhasználói vagy üzleti követelmények megjelenése;
  \item jogszabályi, technológiai vagy környezeti változás;
  \item teljesítmény, biztonság, használhatóság vagy karbantarthatóság javítása;
  \item elavult komponensek, platformok vagy interfészek lecserélése.
\end{itemize}

A karbantartás költsége nagy rendszereknél a fejlesztési költségek többszörösét is elérheti. Ezért a jó architektúra, az olvasható kód, a dokumentáció, a tesztelhetőség és a verziókezelés nem csak fejlesztési kényelmi szempont, hanem hosszú távú gazdasági kérdés is.

\section{Szoftverfolyamat-modellek}

\subsection{A folyamatmodell fogalma}

A \textbf{szoftverfolyamat-modell} a szoftverfolyamat absztrakt reprezentációja. Nem a teljes valóságot írja le, hanem egy bizonyos nézőpontból mutatja meg, hogyan szervezzük a fejlesztési tevékenységeket. A modell segít megérteni, mikor mit kell csinálni, milyen eredményekre épülünk, mikor van visszacsatolás, és hogyan mérhető a projekt előrehaladása.

\begin{center}
\begin{tabularx}{\textwidth}{L{29mm}X X}
\toprule
\textbf{Modell} & \textbf{Alapgondolat} & \textbf{Mikor kedvező?} \\
\midrule
Vízesésmodell & A fő tevékenységeket különálló, egymást követő fázisokként kezeli. Egy fázis eredménye dokumentum, amelyre a következő fázis épít. & Ha a követelmények előre jól ismertek, stabilak, és a projekt jól dokumentált, menedzselhető folyamatot igényel. \\
Evolúciós / iteratív fejlesztés & A specifikáció, a fejlesztés és a validáció tevékenységei összefonódnak; a rendszer több verzión keresztül alakul. & Ha a megrendelő igényei nem teljesen tiszták, gyors visszajelzés kell, vagy a rendszer felhasználói tapasztalat alapján formálódik. \\
Komponensalapú fejlesztés & A rendszer lehetőség szerint meglévő, újrafelhasználható komponensekből épül fel; a követelmények részben igazodhatnak az elérhető komponensekhez. & Ha sok megbízható komponens áll rendelkezésre, és fontos a költség, a kockázat vagy a fejlesztési idő csökkentése. \\
Inkrementális fejlesztés & A rendszer szolgáltatásait kisebb, priorizált inkremensekben szállítják le, miközben a rendszer fokozatosan bővül. & Ha a megrendelőnek korán használható részrendszerre van szüksége, és a funkciók jól priorizálhatók. \\
Spirális fejlesztés & A fejlesztést ismétlődő ciklusokként ábrázolja, amelyekben hangsúlyos a célok kijelölése, a kockázatelemzés, a fejlesztés-validálás és a következő ciklus tervezése. & Nagyobb, kockázatosabb rendszereknél, ahol a kockázatok tudatos kezelése különösen fontos. \\
\bottomrule
\end{tabularx}

\end{center}

A gyakorlatban a modellek ritkán jelennek meg tiszta formában. Egy projekt lehet részben vízesés jellegű a szerződéses dokumentáció miatt, közben alkalmazhat inkrementális szállítást, prototípusokat és komponensalapú fejlesztést is.

\subsection{A vízesésmodell rövid áttekintése}

A \textbf{vízesésmodell} a szoftverfejlesztés egyik klasszikus, korai modellje. A fejlesztést egymás után következő fázisokra bontja: követelmények elemzése, rendszer- és szoftvertervezés, implementáció és egységteszt, integráció és rendszerteszt, majd működtetés és karbantartás.

A modell előnye, hogy egyszerűen áttekinthető, jól dokumentálható és menedzselhető. Minden fázisnak megvan a maga eredménye, például specifikáció vagy tervdokumentum. Ez különösen ott hasznos, ahol szerződéses, biztonsági vagy szabályozási okok miatt részletes dokumentáció kell.

A fő hátránya, hogy a valós fejlesztés ritkán ilyen tisztán lineáris. A követelmények változhatnak, a tervezés közben új kérdések merülhetnek fel, az implementáció során kiderülhet, hogy egy követelmény pontatlan volt. Ha a modell túl merev, akkor a problémák későn derülnek ki, a visszalépések költségesek, és a végén olyan rendszer születhet, amely formálisan megfelel a dokumentumnak, de nem felel meg a valódi igényeknek.

\section{Az evolúciós modell}

\subsection{Az evolúciós fejlesztés alapgondolata}

Az \textbf{evolúciós fejlesztés} alapgondolata az, hogy a fejlesztőcsapat nem próbálja a teljes rendszert előre, véglegesen specifikálni, majd egyetlen nagy lépésben elkészíteni. Ehelyett létrehoz egy kezdeti implementációt, azt megmutatja a felhasználóknak vagy megrendelőnek, visszajelzést gyűjt, majd több verzión keresztül addig finomítja a rendszert, amíg az egyre jobban meg nem felel az igényeknek.

\begin{figure}[htbp]
\centering
\resizebox{0.92\textwidth}{!}{%
\begin{tikzpicture}[font=\scriptsize, node distance=8mm]
  \tikzset{evostep/.style={draw, rounded corners, align=center, minimum width=36mm, minimum height=10mm},
           version/.style={draw, circle, align=center, minimum size=13mm, fill=lightaccent}}

  \node[evostep, fill=warnbg] (known) {Ismert követelmények\\és kezdeti elképzelés};
  \node[evostep, fill=lightaccent, right=15mm of known] (initial) {Kezdeti\\implementáció};
  \node[evostep, fill=black!5, right=15mm of initial] (feedback) {Felhasználói\\véleményezés};
  \node[evostep, fill=lightaccent, right=15mm of feedback] (refine) {Finomítás, javítás,\\új funkciók};
  \node[evostep, fill=warnbg, right=15mm of refine] (target) {Megfelelő\\rendszer};

  \draw[arrow] (known) -- (initial);
  \draw[arrow] (initial) -- (feedback);
  \draw[arrow] (feedback) -- (refine);
  \draw[arrow] (refine) -- (target);
  \draw[arrow, dashed] ([yshift=-2mm]refine.south) .. controls +(0,-18mm) and +(0,-18mm) .. node[below, align=center]{új verzió, újabb visszajelzés} ([yshift=-2mm]initial.south);

  \node[version, below=13mm of initial] (v1) {v1};
  \node[version, below=13mm of feedback] (v2) {v2};
  \node[version, below=13mm of refine] (v3) {v3};
  \draw[arrow] (v1) -- (v2);
  \draw[arrow] (v2) -- (v3);

  \node[draw, rounded corners, fill=black!5, below=10mm of v2, align=center, minimum width=108mm] {Az evolúciós modellben a specifikáció, a fejlesztés és a validáció részben összemosódik: a rendszer a verziók során közelít a valódi igényekhez.};
\end{tikzpicture}%
}
\caption{Az evolúciós fejlesztés alapgondolata}
\label{fig:zv25_evolucios_modell}
\end{figure}


Ebben a modellben a specifikáció, a fejlesztés és a validáció nem különül el olyan élesen, mint a vízesésmodellben. Egy működő részrendszer kipróbálása közben a felhasználó pontosabban meg tudja fogalmazni, mire van szüksége. A fejlesztők közben jobban megértik a problémát, és a rendszer tényleges használata is új követelményeket hozhat felszínre.

Az evolúciós modell különösen hasznos, ha a kezdeti követelmények bizonytalanok, a felhasználói felület vagy üzleti folyamat csak kipróbálás közben értékelhető jól, vagy a megrendelő nem tud mindent előre pontosan megfogalmazni.

\subsection{Az evolúciós fejlesztés két típusa}

\begin{center}
\begin{tabularx}{\textwidth}{L{33mm}X X}
\toprule
\textbf{Típus} & \textbf{Célja} & \textbf{Jellemző működése} \\
\midrule
Feltáró fejlesztés & A megrendelővel együtt fokozatosan feltárni a követelményeket és kialakítani a végleges rendszert. & A fejlesztés a jobban ismert részekkel indul, majd a felhasználói visszajelzések alapján újabb tulajdonságok kerülnek a rendszerbe. \\
Eldobható prototípuskészítés & A bizonytalan, nehezen érthető követelmények pontosítása, nem feltétlenül végleges kód előállítása. & Gyors prototípus készül a kockázatos vagy homályos részekről; a tanulságokat beépítik a specifikációba, a prototípust pedig eldobhatják. \\
\bottomrule
\end{tabularx}

\end{center}

A \textbf{feltáró fejlesztés} során a cél ténylegesen a végleges rendszer fokozatos kialakítása. A csapat a biztosabban ismert részekkel kezd, majd a visszajelzések alapján bővíti és módosítja a rendszert. Itt a létrejövő verziók egyre inkább a végleges termék felé haladnak.

Az \textbf{eldobható prototípuskészítés} célja más. Ilyenkor a prototípus nem feltétlenül a végleges rendszer alapja, hanem tanulási eszköz. Akkor hasznos, ha bizonyos követelmények homályosak, kockázatosak vagy nehezen kommunikálhatók. A prototípus megmutatja, hogy egy elképzelés működik-e, milyen felhasználói reakciókat vált ki, és milyen pontosításokat kell beépíteni a specifikációba.

\subsection{Előnyök és hátrányok}

\begin{center}
\begin{tabularx}{\textwidth}{L{24mm}X}
\toprule
\textbf{Szempont} & \textbf{Értékelés} \\
\midrule
Előny & Gyors visszacsatolást ad, jobban kezeli a bizonytalan követelményeket, és a rendszer a felhasználói tapasztalatok alapján finomítható. \\
Előny & A specifikáció inkrementálisan fejlődhet, ezért az ügyfél hamarabb láthat működő részeket, mint egy szigorúan szekvenciális modellben. \\
Hátrány & A folyamat kevésbé átlátható a menedzsment számára, mert nem mindig vannak jól elkülöníthető, lezárt dokumentum-alapú mérföldkövek. \\
Hátrány & A sok módosítás ronthatja a rendszer szerkezetét, ha nincs tudatos refaktorálás, architekturális kontroll és dokumentálás. \\
Tipikus alkalmazás & Rövidebb élettartamú, kis és közepes rendszereknél, felhasználói felületeknél, webes rendszereknél, illetve olyan projektekben, ahol a követelmények csak használat közben tisztulnak. \\
\bottomrule
\end{tabularx}

\end{center}

Az evolúciós modell legnagyobb előnye a gyors visszacsatolás. A felhasználó nem csak dokumentumot lát, hanem működő rendszerrészletet. Emiatt könnyebben kiderül, hogy a fejlesztők és a megrendelő ugyanazt értik-e egy követelmény alatt.

A modell másik előnye, hogy jobban támogatja a változó igényeket. A rendszer nem egyszeri, lezárt termékként jelenik meg, hanem olyan változatok sorozataként, amelyek egyre közelebb kerülnek a valódi problémához. Ez jól illeszkedik sok mai szoftverhez, ahol az üzleti környezet, a felhasználói elvárás és a technológia is gyorsan változik.

A hátrányok főleg a kontroll és a struktúra területén jelentkeznek. Ha a változtatások nincsenek tudatosan kezelve, a rendszer belső szerkezete romolhat. A gyors javítások, ideiglenes megoldások és folyamatos módosítások technikai adósságot okozhatnak. A menedzsment számára is nehezebb lehet mérni az előrehaladást, mert nem mindig vannak tisztán elkülönülő, lezárt fázisok és dokumentumok.

Ezért az evolúciós fejlesztés nem azt jelenti, hogy nincs tervezés vagy dokumentáció. Inkább azt jelenti, hogy a tervezést, specifikációt és validációt rugalmasabban, iterációkhoz kötve végezzük. Nagyobb rendszereknél különösen fontos az architektúra védelme, a refaktorálás, az automatizált tesztelés és a változások követése.

\subsection{Kapcsolat az inkrementális és spirális modellel}

Az \textbf{inkrementális fejlesztés} köztes szemlélet a vízesés és az evolúciós megközelítés között. A rendszer szolgáltatásait kisebb részekre, inkremensekre bontja. A magasabb prioritású szolgáltatások hamarabb készülnek el, ezért a megrendelő nem csak a teljes rendszer végén kap használható eredményt. Az inkrementumok használata tapasztalatot ad, amely későbbi fejlesztési döntéseket is befolyásolhat.

A \textbf{spirális modell} az iterációt kockázatkezeléssel kapcsolja össze. A fejlesztést ciklusok sorozataként ábrázolja. Egy ciklusban célokat jelölünk ki, azonosítjuk és elemezzük a kockázatokat, választunk egy megfelelő fejlesztési megközelítést, fejlesztünk és validálunk, majd megtervezzük a következő ciklust. A spirális modell fontos sajátossága, hogy expliciten számol a kockázati tényezőkkel, például határidő-, költség-, technológiai vagy követelménykockázatokkal.

Az evolúciós, inkrementális és spirális szemlélet közös pontja, hogy a szoftverfejlesztést nem egyszeri, teljesen előre meghatározott lépéssorként kezeli. Mindegyik elfogadja, hogy a rendszer fokozatosan alakul, a visszacsatolás értékes, és a változás a szoftverfejlesztés természetes része.

\section{Gyakori vizsgahibák}

\begin{warningbox}
\begin{itemize}
  \item \textbf{A szoftvertechnológia leszűkítése programozásra}: a szoftvertechnológia követelménykezelést, tervezést, dokumentációt, validációt, projektvezetést, minőségbiztosítást és karbantartást is jelent.
  \item \textbf{A szoftverfolyamat és a folyamatmodell összekeverése}: a folyamat a tényleges tevékenységek rendszere, a modell ennek absztrakt, egyszerűsített leírása.
  \item \textbf{A fázisok merev sorrendként kezelése}: a specifikáció, tervezés, implementáció, validáció és evolúció a gyakorlatban sokszor visszacsatolásokkal és iterációkkal kapcsolódik össze.
  \item \textbf{Validáció és verifikáció felcserélése}: verifikáció esetén azt nézzük, hogy a specifikációnak megfelelően készül-e a rendszer; validáció esetén azt, hogy a rendszer a valódi célra megfelel-e.
  \item \textbf{Az evolúciós modell azonosítása tervezetlenséggel}: az evolúciós fejlesztésnek is kell architektúra, verziókezelés, tesztelés, dokumentáció és változáskezelés.
  \item \textbf{A prototípus szerepének félreértése}: az eldobható prototípus célja a tanulás és követelménytisztázás, nem feltétlenül végleges kód előállítása.
  \item \textbf{A karbantartás alábecslése}: a szoftver életciklusának jelentős része az evolúció és karbantartás, sokszor ez adja a nagyobb költségrészt.
\end{itemize}
\end{warningbox}

\section{Rövid összefoglaló válasz vizsgára}

\begin{summarybox}
A szoftvertechnológia mérnöki és vezetési alapelvek összessége, amely a szoftverek termékszerű előállítását, dokumentálását, használatát és karbantartását támogatja költség- és határidőkorlátok mellett. A szoftverfolyamat azoknak a tevékenységeknek és eredményeknek a sora, amelyek egy szoftvertermék létrehozásához vezetnek. A közös fázisok: szoftverspecifikáció, tervezés, implementáció, validáció és evolúció. A specifikáció a szolgáltatásokat, megszorításokat és követelményeket rögzíti; a tervezés kialakítja az architektúrát, alrendszereket, interfészeket, adatszerkezeteket és algoritmusokat; az implementáció futtatható rendszerré alakítja a tervet; a validáció és verifikáció ellenőrzi, hogy a rendszer megfelel-e a specifikációnak és a megrendelő céljainak; az evolúció pedig a rendszer életciklusa alatti módosításokat, hibajavításokat és továbbfejlesztéseket jelenti. A szoftverfolyamat-modellek a folyamat absztrakt leírásai. A vízesésmodell különálló, egymásra épülő fázisokat használ, míg az evolúciós modell működő kezdeti implementációból indul ki, felhasználói visszajelzésekkel és több verzión keresztül finomítja a rendszert. Az evolúciós modell előnye a gyors visszacsatolás és a bizonytalan követelmények kezelése, hátránya viszont a nehezebb menedzselhetőség és a rendszerstruktúra romlásának veszélye, ha nincs tudatos tervezés, refaktorálás és változáskezelés.
\end{summarybox}
